% GENERATED FILE. Edit canonical lessons, then run npm run generate:books.
% canonical-source-hash: fnv1a64:d2a5969de2c24f38
% canonical-lessons: KA-C366-khulase, KA-C366-kalanka, KA-C366-ghanate, KA-C366-ane, KA-C366-goladu

\chapter{Disgrace and Dignity}
\label{ch:ka-disgrace-and-dignity}
\hlchaptermodality{366}

\begin{chapteropening}
\textbf{By the end of this chapter:} \emph{I can say an acquittal, a stain, dignity, an oath and to wail.}

Complete the last lesson of chapter 366: I can say an acquittal, a stain, dignity, an oath and to wail.
\end{chapteropening}

\section[\texorpdfstring{\kn{ಖುಲಾಸೆ} (khulāse)}{ಖುಲಾಸೆ (khulāse)}]{\kn{ಖುಲಾಸೆ} (khulāse) — an acquittal}

\label{lesson:KA-C366-khulase}



\begin{quote}
Before the new one: say the Kannada for to shove, then the Kannada for a patrol.
\end{quote}

\subsection*{You'll want to know: \kn{ಖುಲಾಸೆ}}
\textbf{\kn{ಖುಲಾಸೆ}} — \emph{khulāse} — \textquotedblleft{}an acquittal\textquotedblright{}.

\begin{quote}
\kn{ಕೊನೆಗೆ} \kn{ನ್ಯಾಯಾಲಯದಿಂದ} \kn{ಅವನಿಗೆ} \kn{ಖುಲಾಸೆ} \kn{ಸಿಕ್ಕಿತು}.
\end{quote}

\begin{quote}
\emph{konege nyāyālayadinda avanige khulāse sikkitu.}
\end{quote}

\begin{quote}
In the end he was acquitted by the court.
\end{quote}

\kn{ಖುಲಾಸೆ} \kn{ಆಗು}, to be acquitted: \kn{ಆರೋಪಿ} \kn{ಖುಲಾಸೆ} \kn{ಆದನು}.

\begin{cousinweb}
\textbf{\kn{ಖುಲಾಸೆ}} (\emph{khulāse}) came into Kannada through Hindustani and is usually traced to Arabic; the exact Arabic source word is not settled.
\end{cousinweb}

\begin{grammarlens}[title={What you've built}]
One more word for telling what happened.
\end{grammarlens}

\subsection*{Your turn}
\begin{itemize}
\raggedright
  \item \emph{Say it:} \emph{khulāse}
  \item \emph{Say it:} \emph{khulāse}, once more
  \item \emph{Say it:} \emph{konege nyāyālayadinda avanige khulāse sikkitu.} — In the end he was acquitted by the court.
  \item \emph{Recall:} say the Kannada for to shove, then the Kannada for a patrol, then say \emph{khulāse} again
\end{itemize}

\begin{tcolorbox}[breakable,colback=teal!4,colframe=teal!35!black,title={Before you move on}]
What does \kn{ಖುಲಾಸೆ} mean? (\textquotedblleft{}An acquittal\textquotedblright{}.) Say it once more.
\end{tcolorbox}
\section[\texorpdfstring{\kn{ಕಳಂಕ} (kaḷaṅka)}{ಕಳಂಕ (kaḷaṅka)}]{\kn{ಕಳಂಕ} (kaḷaṅka) — a stain; a disgrace}

\label{lesson:KA-C366-kalanka}



\begin{quote}
Before the new one: say the Kannada for a patrol, then the Kannada for an acquittal.
\end{quote}

\subsection*{You'll want to know: \kn{ಕಳಂಕ}}
\textbf{\kn{ಕಳಂಕ}} — \emph{kaḷaṅka} — \textquotedblleft{}a stain; a disgrace\textquotedblright{}.

\begin{quote}
\kn{ಆ} \kn{ಮೋಸ} \kn{ಅವರ} \kn{ಹೆಸರಿಗೆ} \kn{ಕಳಂಕ} \kn{ತಂದಿತು}.
\end{quote}

\begin{quote}
\emph{ā mōsa avara hesarige kaḷaṅka tanditu.}
\end{quote}

\begin{quote}
That fraud brought a stain on his name.
\end{quote}

A stain on a name or a family: \kn{ಕಳಂಕ} \kn{ತರು}.

\begin{cousinweb}
\textbf{\kn{ಕಳಂಕ}} (\emph{kaḷaṅka}) is borrowed from Sanskrit \emph{kalaṅka}, \textquotedblleft{}a stain, a spot\textquotedblright{} (CDIAL 2911), and also \textquotedblleft{}defamation\textquotedblright{} (Monier-Williams).
\end{cousinweb}

\begin{grammarlens}[title={What you've built}]
One more word for telling what happened.
\end{grammarlens}

\subsection*{Your turn}
\begin{itemize}
\raggedright
  \item \emph{Say it:} \emph{kaḷaṅka}
  \item \emph{Say it:} \emph{kaḷaṅka}, once more
  \item \emph{Say it:} \emph{ā mōsa avara hesarige kaḷaṅka tanditu.} — That fraud brought a stain on his name.
  \item \emph{Recall:} say the Kannada for a patrol, then the Kannada for an acquittal, then say \emph{kaḷaṅka} again
\end{itemize}

\begin{tcolorbox}[breakable,colback=teal!4,colframe=teal!35!black,title={Before you move on}]
What does \kn{ಕಳಂಕ} mean? (\textquotedblleft{}A stain; a disgrace\textquotedblright{}.) Say it once more.
\end{tcolorbox}
\section[\texorpdfstring{\kn{ಘನತೆ} (ghanate)}{ಘನತೆ (ghanate)}]{\kn{ಘನತೆ} (ghanate) — dignity}

\label{lesson:KA-C366-ghanate}



\begin{quote}
Before the new one: say the Kannada for an acquittal, then the Kannada for a stain.
\end{quote}

\subsection*{You'll want to know: \kn{ಘನತೆ}}
\textbf{\kn{ಘನತೆ}} — \emph{ghanate} — \textquotedblleft{}dignity\textquotedblright{}.

\begin{quote}
\kn{ಅಜ್ಜ} \kn{ಕಷ್ಟದಲ್ಲೂ} \kn{ಘನತೆಯಿಂದ} \kn{ಬದುಕಿದರು}.
\end{quote}

\begin{quote}
\emph{ajja kaṣṭadallū ghanateyinda badukidaru.}
\end{quote}

\begin{quote}
Grandfather lived with dignity even in hardship.
\end{quote}

\kn{ಘನ}, weighty, with -\kn{ತೆ}, -ness.

\begin{cousinweb}
\textbf{\kn{ಘನತೆ}} (\emph{ghanate}) is built on Sanskrit \emph{ghana}, \textquotedblleft{}compact, firm, dense\textquotedblright{} (CDIAL 4424), with the Kannada abstract ending -\emph{te}: weight of character.
\end{cousinweb}

\begin{grammarlens}[title={What you've built}]
One more word for telling what happened.
\end{grammarlens}

\subsection*{Your turn}
\begin{itemize}
\raggedright
  \item \emph{Say it:} \emph{ghanate}
  \item \emph{Say it:} \emph{ghanate}, once more
  \item \emph{Say it:} \emph{ajja kaṣṭadallū ghanateyinda badukidaru.} — Grandfather lived with dignity even in hardship.
  \item \emph{Recall:} say the Kannada for an acquittal, then the Kannada for a stain, then say \emph{ghanate} again
\end{itemize}

\begin{tcolorbox}[breakable,colback=teal!4,colframe=teal!35!black,title={Before you move on}]
What does \kn{ಘನತೆ} mean? (\textquotedblleft{}Dignity\textquotedblright{}.) Say it once more.
\end{tcolorbox}
\section[\texorpdfstring{\kn{ಆಣೆ} (āṇe)}{ಆಣೆ (āṇe)}]{\kn{ಆಣೆ} (āṇe) — an oath; a vow}

\label{lesson:KA-C366-ane}



\begin{quote}
Before the new one: say the Kannada for a stain, then the Kannada for dignity.
\end{quote}

\subsection*{You'll want to know: \kn{ಆಣೆ}}
\textbf{\kn{ಆಣೆ}} — \emph{āṇe} — \textquotedblleft{}an oath; a vow\textquotedblright{}.

\begin{quote}
\kn{ನಾನು} \kn{ಸತ್ಯ} \kn{ಹೇಳುತ್ತೇನೆ} \kn{ಅಂತ} \kn{ಅವನು} \kn{ಆಣೆ} \kn{ಮಾಡಿದನು}.
\end{quote}

\begin{quote}
\emph{nānu satya hēḷuttēne anta avanu āṇe māḍidanu.}
\end{quote}

\begin{quote}
He swore an oath that he would tell the truth.
\end{quote}

\kn{ಆಣೆ} \kn{ಮಾಡು} or \kn{ಆಣೆ} \kn{ಇಡು}, to swear. \kn{ಆಣೆ} is also an old coin, a sixteenth of a rupee.

\begin{cousinweb}
\textbf{\kn{ಆಣೆ}} (\emph{āṇe}) is usually traced to Sanskrit \emph{ājñā}, \textquotedblleft{}a command\textquotedblright{} (CDIAL 1095), through a Prakrit form \emph{āṇā}.
\end{cousinweb}

\begin{grammarlens}[title={What you've built}]
One more word for telling what happened.
\end{grammarlens}

\subsection*{Your turn}
\begin{itemize}
\raggedright
  \item \emph{Say it:} \emph{āṇe}
  \item \emph{Say it:} \emph{āṇe}, once more
  \item \emph{Say it:} \emph{nānu satya hēḷuttēne anta avanu āṇe māḍidanu.} — He swore an oath that he would tell the truth.
  \item \emph{Recall:} say the Kannada for a stain, then the Kannada for dignity, then say \emph{āṇe} again
\end{itemize}

\begin{tcolorbox}[breakable,colback=teal!4,colframe=teal!35!black,title={Before you move on}]
What does \kn{ಆಣೆ} mean? (\textquotedblleft{}An oath; a vow\textquotedblright{}.) Say it once more.
\end{tcolorbox}
\section[\texorpdfstring{\kn{ಗೋಳಾಡು} (gōḷāḍu)}{ಗೋಳಾಡು (gōḷāḍu)}]{\kn{ಗೋಳಾಡು} (gōḷāḍu) — to wail; to lament}

\label{lesson:KA-C366-goladu}



\begin{quote}
Before the new one: say the Kannada for dignity, then the Kannada for an oath.
\end{quote}

\subsection*{You'll want to know: \kn{ಗೋಳಾಡು}}
\textbf{\kn{ಗೋಳಾಡು}} — \emph{gōḷāḍu} — \textquotedblleft{}to wail; to lament\textquotedblright{}.

\begin{quote}
\kn{ಮನೆ} \kn{ಕಳೆದುಕೊಂಡ} \kn{ಹೆಂಗಸು} \kn{ಗೋಳಾಡಿದಳು}.
\end{quote}

\begin{quote}
\emph{mane kaḷedukoṇḍa heṅgasu gōḷāḍidaḷu.}
\end{quote}

\begin{quote}
The woman who had lost her house wailed.
\end{quote}

Past (I …): \kn{ಗೋಳಾಡಿದೆ}, \emph{gōḷāḍide}.

\kn{ಗೋಳು}, wailing, and \kn{ಆಡು}: to keep on wailing. \kn{ಗೋಳು} \kn{ಹೊಯ್ದುಕೊ} is the same.

\begin{cousinweb}
\textbf{\kn{ಗೋಳಾಡು}} (\emph{gōḷāḍu}) is a native Dravidian compound: \emph{gōḷ}, \textquotedblleft{}the sound of howling, lamenting, wailing\textquotedblright{} (DEDR 2252), and \emph{āḍu}, \textquotedblleft{}to move, to act\textquotedblright{} (DEDR 347).
\end{cousinweb}

\begin{grammarlens}[title={What you've built}]
One more word for telling what happened.
\end{grammarlens}

\subsection*{Your turn}
\begin{itemize}
\raggedright
  \item \emph{Say it:} \emph{gōḷāḍu}
  \item \emph{Say it:} \emph{gōḷāḍu}, once more
  \item \emph{Say it:} \emph{mane kaḷedukoṇḍa heṅgasu gōḷāḍidaḷu.} — The woman who had lost her house wailed.
  \item \emph{Recall:} say the Kannada for dignity, then the Kannada for an oath, then say \emph{gōḷāḍu} again
\end{itemize}

\begin{tcolorbox}[breakable,colback=teal!4,colframe=teal!35!black,title={Before you move on}]
What does \kn{ಗೋಳಾಡು} mean? (\textquotedblleft{}To wail; to lament\textquotedblright{}.) Say it once more.
\end{tcolorbox}
